% Source: https://artofproblemsolving.com/community/c3h2526835_spherical_geo_help
% Stable identity: aops-topic-2526835-q01
In spherical geometry on a sphere with centre $O$, the ``lines'' are the great circles (the circles on the sphere whose centre is $O$), and the angle between two such lines at a point where they cross is defined to be the angle between the planes of the two great circles. Let $N$ be the north pole and let $A$ and $B$ be points on the equator.
\begin{enumerate}
\item[(a)] Show that the angle at $N$ between the great circles through $N$ and $A$ and through $N$ and $B$ equals $\angle AOB$, and that it is also the angle between the tangent lines to these two circles at $N$. (So the definition gives the angle you would see between the two lines when standing at $N$.)
\item[(b)] Suppose $\angle AOB=90^\circ$. Find the three angles of the spherical triangle $NAB$, whose sides are the arcs $NA$, $NB$ and $AB$ of these great circles and of the equator, and find their sum.
\end{enumerate}
